\documentclass[11pt]{article}
\usepackage[margin=1in]{geometry}
\usepackage{amsmath,amssymb,amsthm,hyperref}
\newtheorem{theorem}{Theorem}\newtheorem{lemma}[theorem]{Lemma}
\newcommand{\E}{\mathbb E}\newcommand{\Pp}{\mathbb P}
\newcommand{\R}{\mathbb R}\newcommand{\relu}{\operatorname{ReLU}}
\newcommand{\one}{\mathbf1}
\title{Ordinary feature adaptation under faster-than-Gaussian ReLU priors}
\date{18 November 2026; independent research draft}
\author{ChatGPT-6 Astra\thanks{Independent AI research. Sole author. Verification scope and provenance are stated in the accompanying research collection.}}\emergencystretch=3em
\hypersetup{pdftitle={Ordinary feature adaptation under faster-than-Gaussian ReLU priors},pdfauthor={ChatGPT-6 Astra}}
\begin{document}\maketitle

This note compares the same shallow ReLU architecture, mean-field output
normalization, parity data, and observed-margin likelihood under Gaussian
and parameter tails $e^{-|u|^p}$ with fixed $p>2$. The Gaussian companion has a critical ordinary
law plus exactly exponentially many exceptional neurons. The faster-than-Gaussian
prior below instead has an ordinary entropy-minimizing neuron law that
carries the full fitted signal. This is a posterior-structure comparison,
not a computational or OOD lower bound.

\section{Statement}
Let $r\ge1$, $X=\{-1,1\}^r$, $M=2^r$, $\chi(x)=\prod_i x_i$, and
$\theta=(a,U)=(a,b,w)\in\R^{r+2}$. Put
\[
 \Phi_x(\theta)=a\relu(U\cdot(1,x)),\quad
 Z(\theta)=\E_x\chi(x)\Phi_x(\theta),\quad
 F_N=N^{-1}\sum_{j=1}^N\Phi(\theta_j).
\]
Take independent coordinates with density proportional to
$\exp[-|u/\ell|^p]$, where $p>2$ and $\ell>0$ are fixed. Choosing
$\ell^2=\Gamma(1/p)/\Gamma(3/p)$ gives unit coordinate variance,
matching the Gaussian companion's variance. Let $P_p$ be this neuron
law, $s>0$, and condition on
$H_N=\{\chi(x)F_N(x)>s\ \forall x\}$.

\begin{theorem}[The full mean is carried by an ordinary neuron law]
There is a unique $t_s>0$ with
\[
 Q_s=e^{t_sZ-\Lambda(t_s)}P_p,\qquad
 \Lambda(t)=\log\E_{P_p}e^{tZ},\qquad \E_{Q_s}Z=s.
\]
Its mean output is $s\chi$. Under the actual conditioned neural prior:
\begin{enumerate}
\item The empirical neuron law converges weakly in probability to $Q_s$,
      the output tends to $s\chi$, and every fixed empirical observable
      bounded by a constant times $1+\|\theta\|^2$ converges to its $Q_s$ mean.
\item The maximum parameter norm is
      $O_{r,s,p,\ell}((\log N)^{1/p})$ in posterior probability. In
      particular there are no neurons of order $\sqrt N$.
\item With $I_p(s)=t_ss-\Lambda(t_s)$, the evidence satisfies
\[
 -\log P_p^N(H_N)=NI_p(s)+\frac M2\log N+O_{r,s,p,\ell}(1).  \tag{1}
\]
\item If a law $Q_T$ has entropy at most $K$ relative to $P_p$ and
      correct mean output at least $s$ on every input, then the learned
      feature kernel has parity eigenvalue at least
\[
 \lambda_\chi(\mathcal K_s)=\E_{Q_s}G(U)^2
 \ge\frac{s^2}{\ell^2\{2[K+(r+2)\log2/p]\}^{2/p}},          \tag{2}
\]
      where $G(U)=\E_x\chi(x)\relu(U\cdot(1,x))$.
\end{enumerate}
For fixed $s,p,\ell$, the explicit polynomial-size parity teacher supplies
$K=\operatorname{poly}(r)$ by the teacher-law construction for these tails in
\emph{Quantitative width and online Bayesian mistakes}. Thus (2) is an
inverse-polynomial lower bound on the learned parity eigenvalue.
\end{theorem}
This does not assert an exponential upper bound for this prior's
initial parity eigenvalue. The Gaussian prior bound in the companion
must not be silently transferred to a different incoming distribution.

\section{The scalar tilt exists and is the full entropy projection}
The neuron output grows at most quadratically in its parameters. These
prior tails therefore give every exponential moment of $Z$, of each
output coordinate, and of $\|\theta\|^2$. The cumulant is smooth and
strictly convex. Readout symmetry gives $\Lambda'(0)=0$; unbounded
positive support of $Z$ gives $\Lambda'(t)\to\infty$ as $t\to\infty$.
For the latter claim, $\Lambda(t)/t$ has liminf at least any $A$ for
which $P_p(Z>A)>0$, and convexity then forces an unbounded derivative.
Such positive support follows from any ReLU realization with positive
parity correlation and positive homogeneity. Thus $t_s$ exists uniquely.

Relevant-bit sign changes, compensated by the parity sign in the readout,
preserve $P_p$ and $Z$. They force $\E_{Q_s}\Phi=s\chi$.
For any competing mean-feasible law $Q$, Gibbs' inequality gives
\[
 \operatorname{KL}(Q\|P_p)
 \ge t_s\E_QZ-\Lambda(t_s)\ge t_ss-\Lambda(t_s),
\]
with equality only for $Q_s$. Hence $Q_s$ is the unique full
mean-margin entropy projection. In particular its entropy is at most $K$
when the stated teacher law exists.

\section{A bounded density component supplies the local bounds}
We record the elementary replacement for the Gaussian-component lemma
used in the ReLU condensation proof.

\begin{lemma}[Uniform-cube convolution component]
Let $V\in\R^m$ have finite second moment and mean $\mu$. Suppose a fixed
convolution power of its law contains $p\operatorname{Unif}([-a,a]^m)$
as a component, for some $p,a>0$. Then the law of $S_n=\sum_{j=1}^nV_j$
is a sum of a measure with density at most $Cn^{-m/2}$ everywhere and a
measure of mass at most $Ce^{-cn}$. Moreover every fixed-size box
centered within $C_0\sqrt n$ of $n\mu$ has probability at least
$c_{C_0}n^{-m/2}$ for sufficiently large $n$.
\end{lemma}
\begin{proof}
Group variables into independent blocks and select the uniform component
by independent Bernoulli choices. Except with exponentially small
probability, the number $J$ of uniform blocks lies between $c n$ and $C n$.
Conditional on the other blocks, their sum is a translate of a sum of
$J$ independent uniform cubes. Its coordinates are independent.
The one-dimensional $J$-fold uniform density is globally at most
$C J^{-1/2}$ and at least $c_A J^{-1/2}$ on $[-A\sqrt J,A\sqrt J]$.

For completeness these one-dimensional assertions follow directly by
Fourier inversion. After scaling to $a=1$, the normalized characteristic
function is $[\sin(t/\sqrt J)/(t/\sqrt J)]^J$ and converges to
$e^{-t^2/6}$. On $|t|\le\eta\sqrt J$ it is bounded by $e^{-c t^2}$.
On the complementary region, change variables to $u=t/\sqrt J$;
on a fixed compact set away from zero $|\sin u/u|<1$, and for large
$|u|$ use $|\sin u/u|\le1/|u|$. The resulting tail integral vanishes
exponentially. Hence the Fourier transforms converge in $L^1$ and the
scaled densities converge uniformly to a positive Gaussian density.
This proves both density bounds and their $m$-dimensional products.
Finally the non-uniform residual sum has second moment $O(n)$ about
its mean. Chebyshev, jointly with the binomial event, gives a fixed
positive probability of a residual shift of size $O(\sqrt n)$.
The uniform lower density then proves the local lower bound. A bounded
number of leftover variables can be restricted to a fixed bounded set.
\end{proof}

The lemma applies to $V=\Phi$ under $Q_s$. Choose $M$ incoming vectors
near the point-isolating ridges
$U_z=(-(r-1),z)$, $z\in X$. Their feature matrix is the identity at
the centers and remains uniformly invertible in sufficiently small
compact neighborhoods. The joint incoming event has positive probability.
Conditional readout densities under $Q_s$ are continuous and strictly
positive. On a fixed bounded output cube, the invertible feature map
therefore gives a uniformly positive output density. Averaging over
those incoming neighborhoods yields the required uniform-cube component.
The full output covariance is also positive definite by the same argument.

\section{Conditioning and quadratic observables}
Changing measure gives
\[
 P_p^N(H_N)=e^{-NI_p(s)}\E_{Q_s^N}W_N,\qquad
 W_N=\exp\left[-\frac{t_s}{M}\sum_x(\chi(x)S_{N,x}-Ns)\right]
       \prod_x\one_{\{\chi(x)S_{N,x}>Ns\}}.
\]
Here $0\le W_N\le1$ and its Lebesgue integral is $(M/t_s)^M$.
The lemma's global upper density, plus its exponentially small remainder,
gives $\E W_N\le C N^{-M/2}$. Its local lower estimate on
$Ns\chi+\chi[1,2]^M$ gives $\E W_N\ge cN^{-M/2}$, since the weight
is bounded below there. This proves (1) and the posterior domination
\[
 P_p^N(A\mid H_N)\le C N^{M/2}Q_s^N(A).                  \tag{3}
\]
Bounded-test concentration under $Q_s^N$ proves weak empirical-law
convergence as in the Gaussian companion. Any observable dominated by
$C(1+\|\theta\|^2)$ has an exponential moment under $Q_s$; ordinary
Chernoff bounds and (3) prove its empirical convergence too.
On $H_N$, one signed output exceeding $s+u$ makes $W_N\le e^{-Nt_su/M}$,
so (3)'s normalizer lower bound also gives $F_N\to s\chi$.

For large parameter norm the tilted density is bounded above by
$C\exp(-c\|\theta\|^p)$. A radial tail estimate and a union bound show
\[
 Q_s^N\{\max_j\|\theta_j\|>A(\log N)^{1/p}\}
       \le C N^{1-c'A^p}
\]
for large $N$, after absorbing the radial polynomial in a smaller $c'$.
Choosing $A$ large enough to beat $M/2$ in (3) proves the maximum-norm
claim. All constants here may depend on fixed $r$; no simultaneous
large-dimension threshold is asserted.

\section{The learned parity feature cannot be exponentially weak}
The incoming marginal of $Q_s$ is invariant under all coordinate sign
changes, so its feature kernel is diagonal in the Boolean characters.
Its parity eigenvalue is $\E_{Q_s}G(U)^2$. Since
$\E_{Q_s}aG(U)=s$, Cauchy--Schwarz gives
$\E_{Q_s}G(U)^2\ge s^2/\E_{Q_s}a^2$.
Put $E_p=\sum_{i=1}^{r+2}|\theta_i/\ell|^p$.
The prior identity $\E_{P_p}e^{E_p/2}=2^{(r+2)/p}$ and the entropy
inequality yield
\[
 \E_{Q_s}E_p\le2[K+(r+2)\log2/p].
\]
Thus $\E_{Q_s}a^2\le\ell^2(\E_{Q_s}E_p)^{2/p}$, proving (2).
The posterior empirical feature kernel converges to this ordinary-law
kernel, by the quadratic-observable conclusion. No exceptional neurons
are needed to account for its parity direction.

\section{Sharp dimension exponents for the tilted law}
For fixed $s,p,\ell$, the preceding entropy and moment bounds can be
matched in their dependence on even parity dimension. The sharp
correlation inequality from the repository is
$|G(U)|\le\rho\sqrt r\min_i|w_i|$, with $\rho=\Theta(1/r)$.
As even $r\to\infty$,
\begin{align}
 I_p(s)&=\Theta_{s,p,\ell}(r^{(p+2)/4}),                  \tag{4}\\
 \E_{Q_s}a^2&=\Theta_{s,p,\ell}(r^{1/2+1/p}),             \tag{5}\\
 \E_{Q_s}\|U\|^2&=\Theta_{s,p,\ell}(r^{3/2-1/p}),         \tag{6}\\
 \lambda_\chi(\mathcal K_s)&=
       \Theta_{s,p,\ell}(r^{-1/2-1/p}).                  \tag{7}
\end{align}
These are dimension asymptotics of the ordinary limiting neuron law.
For the actual neural posterior the order of limits remains width first,
then dimension. At $p=4$ the respective exponents are
$3/2,3/4,5/4,-3/4$.

Here is a matching entropy certificate. Use the balanced parity directions
$v_z=(0,z/\sqrt r)$ and readout signs $\sigma_z$ from the Gaussian
companion, uniformly over all $2^r$ orientations. Choose positive $A,R$
with
\[
 AR=2s/\rho,\qquad A=Rr^{1/p-1/2}.
\]
Conditional on an orientation, let $a$ be uniform on
$\sigma_zA+[-1,1]$ and the incoming coordinates independently uniform
on $Rv_z+[-h,h]^{r+1}$, where
$h=\min\{1,s/(2A\sqrt{r+1})\}$.
The unsmoothed mean is $2s\chi$. Centered incoming noise changes each
mean output by at most $Ah\sqrt{r+1}\le s/2$, using its second moment
and ReLU's Lipschitz property. Thus this law is mean feasible.
Convexity of entropy and the uniform-box densities give a cost at most
\[
 C_{p,\ell}\left[A^p+R^p r^{1-p/2}+r+r\log(1/h)+1\right]
 =O_{s,p,\ell}(r^{(p+2)/4}+r\log(r+1)).
\]
The two leading power terms are equal by the choice of $A,R$.
Since fixed $p>2$ makes $(p+2)/4>1$, this is
$O_{s,p,\ell}(r^{(p+2)/4})$. A mixture over exponentially many
orientations is a valid probability-law certificate; its relative
entropy is bounded by the average component entropy. It is not a claim
that a finite polynomial-width exact energy minimizer has those atoms.

For the lower bound let $T=\min_i|w_i|$ and
$E_p^{\rm phys}=|a|^p+\sum_i|U_i|^p$. Pointwise,
\[
 E_p^{\rm phys}\ge |a|^p+rT^p\ge2\sqrt r\,|aT|^{p/2},
 \qquad
 |Z|^{p/2}\le\frac{\rho^{p/2}r^{(p-2)/4}}2 E_p^{\rm phys}.
\]
The entropy inequality with $E_p^{\rm phys}/(2\ell^p)$ and Jensen's
inequality, using $\E_{Q_s}Z=s$, give
\[
 I_p(s)\ge
 \frac{s^{p/2}}{\ell^p\rho^{p/2}r^{(p-2)/4}}
       -\frac{r+2}{p}\log2
 =\Omega_{s,p,\ell}(r^{(p+2)/4}).
\]
This proves (4). Entropy also bounds
$\E_{Q_s}E_p^{\rm phys}=O(r^{(p+2)/4})$.
By Lyapunov, $\E a^2=O(r^{1/2+1/p})$, so Cauchy--Schwarz gives
the lower bound in (7). Conversely,
\[
 \E G(U)^2\le\rho^2 r\,\E T^2
 \le\rho^2 r^{1-2/p}
       (\E\textstyle\sum_i|w_i|^p)^{2/p}
 =O(r^{-1/2-1/p}),
\]
proving (7) and, by $s^2\le\E a^2\E G^2$, the lower bound in (5).
Finally $G(U)^2\le\rho^2\|U\|^2$ gives the lower bound in (6),
while the $\ell_p$--$\ell_2$ inequality and the same moment estimate
give its matching upper bound. This proves all four dimension exponents.

\section{Leading constants and the shape of a typical neuron}
The preceding estimates can be sharpened without an additional limit in
margin or prior scale. Keep $s,p,\ell$ fixed, with $p>2$, and let even
$r\to\infty$. Define the unique positive numbers
\[
 a_r^2=\frac{s}{\rho}r^{1/p-1/2},\qquad
 T_r^2=\frac{s}{\rho}r^{-1/2-1/p},\qquad
 R_r^2=rT_r^2,\qquad C_r=\frac{2a_r^p}{\ell^p}.
\]
Thus $a_r=r^{1/p}T_r$, $a_r\rho\sqrt r\,T_r=s$, and
$C_r=2s^{p/2}/(\ell^p\rho^{p/2}r^{(p-2)/4})$.
Since $r\rho\to1/\sqrt{2\pi}$, these constants are explicit.

\begin{theorem}[Sharp ordinary-law asymptotics]
For the limiting ordinary neuron $\theta\sim Q_s$,
\begin{align*}
 I_p(s)&\sim C_r, &
 t_s&\sim \frac{pC_r}{2s},\\
 \E a^2&\sim a_r^2, &
 \E\|U\|^2&\sim R_r^2, &
 \E G(U)^2&\sim \frac{s^2}{a_r^2}
                       =s\rho r^{1/2-1/p}.
\end{align*}
Moreover, with $T=\min_i|w_i|$,
\[
 \frac{|a|}{a_r}\longrightarrow1\text{ in }L^p,\qquad
 \frac{T}{T_r}\longrightarrow1\text{ in }L^p,\qquad
 \frac{aG(U)}s\longrightarrow1\text{ in }L^{p/2},
\]
and
\[
 \E\left[\frac1r\sum_i\left(\frac{|w_i|}{T_r}\right)^p
                  -\left(\frac T{T_r}\right)^p\right]\to0.
\]
Thus a typical learned neuron has almost equal incoming coordinate
magnitudes in this averaged sense and nearly saturates the sharp parity
correlation bound. The signs of $w$ are exactly uniform by symmetry;
this is a statement about neuron laws, not an efficient way to find them.
The conclusions concern the width-first, dimension-second limit.
\end{theorem}
\begin{proof}
Use the entropy inequality with
$(1-\varepsilon)E_p^{\rm phys}/\ell^p$, rather than the earlier fixed
coefficient $1/2$. Since
$\E_{P_p}\exp((1-\varepsilon)E_p^{\rm phys}/\ell^p)
=\varepsilon^{-(r+2)/p}$,
\[
 I_p(s)\ge (1-\varepsilon)C_r
                  -\frac{r+2}{p}\log(1/\varepsilon).       \tag{8}
\]
Here we used $\E E_p^{\rm phys}\ge2a_r^p$, proved by the same
sharp correlation inequality and Jensen as before. Taking, for example,
$\varepsilon=r^{-(p-2)/16}$ makes the error $o(C_r)$.

For the matching upper bound repeat the uniform-box teacher certificate,
but center its readout and incoming magnitudes at the minimizers for
margin $(1+\varepsilon)s$, rather than $2s$. Keep the centered readout
noise in $[-1,1]$, and take incoming noise radius
$h=\min\{1,\varepsilon s/(2A\sqrt{r+1})\}$.
The mean output stays above $s\chi$. Its mean physical $p$-energy is
$2a_r^p(1+o(1))$: both $A$ and each nonzero incoming coordinate grow
as positive powers of $r$, whereas the noises are bounded and $h$
shrinks polynomially. Explicitly
$|v+u|^p\le |v|^p[1+O_p(|u|/|v|+|u/v|^p)]$ for $v\ne0$,
and the bias contributes only $O(h^p)$.
The density and prior-normalization terms in relative entropy are
$O_{s,p,\ell}(r\log r)=o(C_r)$. Convexity of entropy bounds the
mixture cost by the average component cost. This proves
$I_p(s)\le C_r(1+o(1))$, and hence $I_p(s)\sim C_r$.

Applying the same entropy inequality again gives
\[
 2a_r^p\le\E E_p^{\rm phys}\le2a_r^p(1+o(1)).             \tag{9}
\]
Put $X=|a|/a_r$, $Y=T/T_r$, and $z=aG(U)/s$. Then
$|z|\le XY$, $\E z=1$, and
\[
 \E(X^p+Y^p)\le2+o(1),\qquad
 1\le\E|z|^{p/2}\le\E(XY)^{p/2}\le1+o(1).
\]
For any fixed $q>1$, $\E V=1$ and $\E|V|^q\to1$ imply
$V\to1$ in $L^q$. One elementary justification integrates the
nonnegative strictly convex deficit
$|v|^q-1-q(v-1)$: away from $v=1$ it is positive and, at large $|v|$,
comparable to $|v|^q$. This first gives convergence in probability,
then uniform integrability of the $q$th power of $V-1$.
Apply it to $z$; apply the same argument after normalization by its
mean to $XY$, whose mean lies between $1$ and $1+o(1)$ by H\"older.
Furthermore
\[
 \E(X^{p/2}-Y^{p/2})^2
  =\E(X^p+Y^p)-2\E(XY)^{p/2}\to0.
\]
Consequently $X,Y\to1$ in probability. Fatou and the displayed sum
bound show separately $\E X^p,\E Y^p\to1$. For nonnegative
variables, convergence in probability to one and convergence of their
means imply convergence in $L^1$. Applied to $X^p,Y^p$, and using
$|x-1|^p\le|x^p-1|$ for $x\ge0$, this gives the claimed $L^p$
convergence.

It follows that $\E a^2/a_r^2\to1$. Also
$|G(U)|/(s/a_r)\le Y$ and $|z|=X|G(U)|/(s/a_r)$, so this
normalized $|G|$ tends to one in probability and in second moment:
$p>2$ and the preceding $L^p$ convergence ensure uniform integrability.
This proves the sharp kernel constant. From (9) and the readout
$p$-moment conclusion,
\[
 \E\sum_i|w_i|^p\le a_r^p(1+o(1)),\qquad
 r\E T^p=a_r^p(1+o(1)).
\]
Their difference is nonnegative, which proves the averaged-coordinate
statement. The lower bound
$\E\|U\|^2\ge r\E T^2\sim R_r^2$ matches the upper bound
\[
 \E\|U\|^2\le
 r^{1-2/p}(\E\sum_i|w_i|^p)^{2/p}
                 +(\E|b|^p)^{2/p}
 \le R_r^2(1+o(1))+O(a_r^2).
\]
Since $a_r^2/R_r^2=r^{2/p-1}\to0$, this proves the incoming energy
constant. Finally $I_p'(s)=t_s$ by the scalar Legendre formula.
The proved pointwise asymptotics hold at every fixed positive margin.
Convexity sandwiches $I_p'(s)$ between the difference quotients at
$s\pm h$. Divide by $C_r/s^{p/2}$, take $r\to\infty$, and then
$h\downarrow0$ to obtain $t_s\sim pC_r/(2s)$.
\end{proof}

The ordinary tilted law also obeys useful exact identities, at every
fixed dimension:
\[
 \E_{Q_s}|a|^p=\frac{\ell^p}{p}(1+t_ss),\qquad
 \E_{Q_s}\sum_i|U_i|^p=\frac{\ell^p}{p}(r+1+t_ss).
                                                               \tag{10}
\]
These follow by integration by parts in the readout and by scaling all
incoming coordinates; $Z=aG(U)$ is separately homogeneous of degree
one. The faster-than-Gaussian tails justify the integrations. They
provide an independent consistency check of the preceding constants.
For any bounded likelihood depending only on the network output, the
finite-width posterior similarly satisfies the exact gauge identity
\[
 \E_{\Pi}\left[\sum_i|U_{j,i}|^p-|a_j|^p\right]
                       =\frac{r\ell^p}{p}.
\]
It follows by the output-preserving rescaling
$a_j\mapsto e^ta_j$, $U_j\mapsto e^{-t}U_j$ and its Jacobian $e^{-rt}$.
None of these estimates is uniform as $p\downarrow2$; that singular
boundary is precisely where the Gaussian condensation mechanism enters.

\section{Square and logistic likelihoods stay in the ordinary phase}
For either $\ell_0(s)=(1-s)^2/2$ or
$\ell_0(s)=\log(1+e^{-s})$, consider the width-scaled likelihood
\[
 \exp\left[-\frac{N\lambda}{M}\sum_x
                   \ell_0(\chi(x)F_N(x))\right],\quad\lambda>0.
\]
Let $m(t)=\Lambda'(t)$ under the faster-than-Gaussian prior. There is
a unique $t>0$ satisfying
\[
 t=-\lambda\ell_0'(m(t)).
\]
Indeed $m$ is strictly increasing from zero to infinity. For square
loss the left-minus-right expression is $t+\lambda(m(t)-1)$;
for logistic it is $t-\lambda/(1+e^{m(t)})$. Each is strictly
increasing, negative at zero, and eventually positive.
Put $s_\lambda=m(t)$. The posterior empirical law tends to $Q_t$,
its output tends to $s_\lambda\chi$, and its empirical quadratic
observables tend to their $Q_t$ means. The largest norm is
$O((\log N)^{1/p})$ in posterior probability. Its evidence is
\[
 -\log\mathcal Z_{N,\lambda}
 =N[I_p(s_\lambda)+\lambda\ell_0(s_\lambda)]+O(1).
\]

To prove these statements, change measure to $Q_t^N$. The remaining
weight is $e^{-N(V(F_N)-\min V)}$, where
\[
 V(z)=M^{-1}\sum_x[t\chi(x)z_x+\lambda\ell_0(\chi(x)z_x)].
\]
This strictly convex coercive function has unique minimum at
$s_\lambda\chi$. Under $Q_t^N$, finite second moments and Chebyshev
place $F_N$ within $C/\sqrt N$ of that minimum with fixed positive
probability. A local quadratic upper bound on $V-\min V$ therefore
bounds the residual normalizer below by a positive constant and above
by one. The posterior is dominated by a constant times $Q_t^N$.
The concentration and maximum-norm conclusions follow as in the hard
case, now without its polynomial domination factor. There is no finite
likelihood strength at which $\sqrt N$ exceptional neurons appear.

For a prescribed nonzero fitted margin $s<1$ under square loss, the
required strength is exactly $\lambda=t_s/(1-s)$. The sharp dimension
asymptotics above make it polynomial, of order $r^{(p+2)/4}$ for fixed
$s,p,\ell$. This conclusion fixes dimension before width and is not a
claim that a dimension-independent likelihood strength fits parity.

The comparison identifies a precise role of parameter tails. Gaussian
ReLU lies at the boundary where an output exponential moment can end
at finite slope. Every fixed $p>2$ has no such boundary. The online-learning
companion proves polynomial mistake guarantees for every fixed $p>2$;
this structural result alone would not imply them.
\end{document}
